\documentclass[10pt,a4paper]{amsart}
\usepackage[margin=19mm]{geometry}
\usepackage{amsmath,amssymb,mathtools}
\usepackage[scr=rsfs,cal=euler]{mathalfa}
\usepackage{xfrac}
\usepackage[unicode,hidelinks,bookmarks=false]{hyperref}
\newcommand{\ie}{i.e.\ }
\newcommand{\cf}{cf.\ }
\newcommand{\resp}{respectively\ }
\newcommand{\Subsection}{\subsection}
\providecommand{\nobreakdash}{\nobreak\hbox{-}}
\setlength{\emergencystretch}{2em}
\multlinegap0pt
\usepackage{bm}
\usepackage{bbm}
\usepackage{dsfont}
\usepackage{xspace}
\usepackage{cancel}
\usepackage{mathtools}
\usepackage{amsthm}
\makeatletter
\@ifundefined{theorem}{\newtheorem{theorem}{Theorem}}{}
\@ifundefined{assumption}{\newtheorem{assumption}[theorem]{Assumption}}{}
\@ifundefined{lemma}{\newtheorem{lemma}[theorem]{Lemma}}{}
\makeatother
\makeatletter
\expandafter\ifx\csname proof\endcsname\relax
\newenvironment{proof}[1][Proof]{\textbf{#1.} }
\fi
\makeatother
\title[$\ell\_{1}^{2}-\eta\ell\_{2}^{2}$ sparsity regularization for ]{$\ell\_{1}^{2}-\eta\ell\_{2}^{2}$ sparsity regularization for nonlinear ill-posed problems}
\author{Long Li, Liang Ding}
\date{Source: arXiv:2508.16163v1 (CC BY 4.0)}
\begin{document}
\maketitle

\noindent\emph{Source and licence.} $\ell_{1}^{2}-\eta\ell_{2}^{2}$ sparsity regularization for nonlinear ill-posed problems. Version arXiv:2508.16163v1 (CC BY 4.0), distributed under the Creative Commons Attribution 4.0 International licence, as recorded in the arXiv licence metadata for this exact version and shown on the abstract page https://arxiv.org/abs/2508.16163v1. Full rights notice: licenses/arxiv-2508.16163-CC-BY-4.0.txt.\par\smallskip

\noindent\textit{Editorial scope.} Counted unit: the Lemma whose source label is \texttt{lem3.4} in the article (source0.tex lines 522--528, source label \texttt{lem3.4}), delivered together with the article's own complete original proof of it (source0.tex lines 530--611, one \texttt{proof} environment of 82 lines). No mathematics is written, repaired or re-derived: statement and proof are the article's own text line for line. Required context retained uncounted: equ1.1 (lines 127--129); equ1.3 (lines 156--158); equ2.2 (lines 194--196); ass3.1 (lines 482--494); equ3.1 (lines 485--487); equ3.4 (lines 504--506); lem3.3 (lines 510--512); equ3.5 (lines 517--520); equ3.6 (lines 517--520). Every definition, display and label of those lines is retained unchanged. Omitted from this package and not redistributed: 1--126, 130--155, 159--193, 197--481, 488--503, 507--509, 513--516, 521--521, 529--529, 612--1444. Nothing inside a retained excerpt is omitted or altered: every retained body is delivered exactly as the article's lines are written, so no omission receipt of any kind accompanies this package. Citations: the retained excerpts carry no citation command at all, so no bibliography entry of the article is transcribed and no reference is redistributed. Reference adaptation: none was needed and none was made; every reference inside the delivered unit resolves inside it, so no unresolved reference or citation is produced. Printed slips kept literal: no printed word, symbol, hypothesis or number of the counted statement or of its proof was altered. Layout and class: the article's own class is \texttt{article} with packages including latexsym, mathrsfs, amsfonts, amsmath, amssymb, indentfirst, subeqnarray, verbatim, graphicx, epstopdf, algorithm, algorithmic, subfigure, color; the wrapper is the lane's pinned \texttt{amsart} editorial wrapper at 10pt on A4, which restates the theorem-like environments the retained text needs and adds the article's own macro definitions that the retained text needs (none), with extra wrapper packages (bm, bbm, dsfont, xspace, cancel, mathtools). The delivered numbering is the unit's own. Assets: no retained span contains a figure environment, a graphics-inclusion command, a PDF-page inclusion, a table or a TikZ picture; no image file is copied into this package, the package declares no asset dependency and no image file is redistributed with it. Licence: version arXiv:2508.16163v1 of this article is distributed under the Creative Commons Attribution 4.0 International licence, as recorded in the arXiv licence metadata for this exact version and shown on the abstract page https://arxiv.org/abs/2508.16163v1. A full rights notice is delivered at licenses/arxiv-2508.16163-CC-BY-4.0.txt. Characters: every retained excerpt is pure ASCII, so the delivered engine, XeLaTeX with its default Latin Modern text font, reports no missing character.
\begin{equation}\label{equ1.1}
     F\left(x\right)=y,
\end{equation}

\begin{equation}\label{equ1.3}
    \mathcal{R}_{\eta}\left(x\right):=\left\|x\right\|_{\ell_{1}}^{2} -\eta\left\|x\right\|_{\ell_{2}}^{2},\quad 0<\eta\leq 1.
\end{equation}

\begin{equation}\label{equ2.2}
    I\left(x^{\dagger}\right)=\left\{i\in\mathbb{N}\mid x_{i}^{\dagger}\neq0\right\},
\end{equation}

\begin{assumption}\label{ass3.1}
    Let $x^{\dagger}\neq 0$ be an $\mathcal{R}_{\eta}$-minimizing solution to the problem \eqref{equ1.1}, which exhibits sparsity. Assume that \\
    (1) $F$ is continuously Fr\'{e}chet-differentiable. For each $i\in I\left(x^{\dagger}\right)$, there exists $\omega_{i}\in Y$ such that
    \begin{equation}\label{equ3.1}
        e_{i} = F'\left(x^{\dagger}\right)^{*}\omega_{i},
    \end{equation}
    where $I\left(x^{\dagger}\right)$ is defined in \eqref{equ2.2}. \\
    (2) There exists constants $\gamma>0$, $\rho>0$, and $\mathcal{R}_{\eta}\left(x^{\dagger}\right)<\rho$ such that
    \begin{equation}\label{equ3.2}
        \left\|F'\left(x\right)-F'\left(x^{\dagger}\right)\right\|_{L\left(\ell_{2}, Y\right)}\leq \gamma\left\|x-x^{\dagger}\right\|_{\ell_{2}}
    \end{equation}
    for all $x\in\ell_{2}$ that satisfy $\mathcal{R}_{\eta}\left(x\right)\leq \rho$.
\end{assumption}

    \begin{equation}\label{equ3.4}
        \left\|F'\left(x^{\dagger}\right)\left(x-x^{\dagger}\right)\right\|_{Y}\leq \frac{\gamma}{2}\left\|x-x^{\dagger}\right\|_{\ell_{2}}^{2}+ \left\|F\left(x\right)-F\left(x^{\dagger}\right)\right\|_{Y}
    \end{equation}

\begin{lemma}\label{lem3.3}
    If Assumption \eqref{equ3.1} holds, it follows that there exists $x^{\dagger}= F'\left(x^{\dagger}\right)^{*}\omega^{\dagger}$. 
\end{lemma}

\begin{align}
    c_{1} &= \left(M_{1}+2\left\|x^{\dagger}\right\|_{\ell_{1}}\right) \left|I\left(x^{\dagger}\right)\right|\max_{i\in I\left(x^{\dagger}\right)} \left\|\omega_{i}\right\|_{Y}+2\left\|\omega^{\dagger}\right\|_{Y},  \label{equ3.5} \\
    c_{2} &= 2\left\|\omega^{\dagger}\right\|_{Y}.  \label{equ3.6}
\end{align}

\begin{lemma}\label{lem3.4}
    Let $M>0$ be given, and denote by $c_{1}$ and $c_{2}$ the constants defined in \eqref{equ3.5}-\eqref{equ3.6}. Under Assumption \ref{ass3.1}, if $\Gamma:=\frac{\gamma \left(c_{1}-c_{2}\eta\right)}{2\left(1-\eta\right)}<1$, then we have
    \begin{equation*}
        \left\|x-x^{\dagger}\right\|_{\ell_{2}}^{2}\leq \frac{1}{1-\Gamma} \left[\frac{1}{1-\eta}\left(\mathcal{R}_{\eta}\left(x\right)- \mathcal{R}_{\eta}\left(x^{\dagger}\right)\right)-\frac{c_{1}-c_{2}\eta} {1-\eta}\left\|F\left(x\right)- F\left(x^{\dagger}\right)\right\|_{Y}\right]
    \end{equation*}
    for any $x\in\ell_{2}$ satisfying $\mathcal{R}_{\eta}\left(x\right)\leq M$ and $\mathcal{R}_{\eta}\left(x\right)\leq \rho$.
\end{lemma}

\begin{proof}
    By the definition of $\mathcal{R}_{\eta}\left(x\right)$ in \eqref{equ1.3}, we can expound that 
    \begin{align*}
        \mathcal{R}_{\eta}\left(x\right)
        & = \left\|x\right\|_{\ell_{1}}^{2}-\eta\left\|x\right\|_{\ell_{2}}^{2} = \left(\left\|x\right\|_{\ell_{1}}^{2}-\left\|x\right\|_{\ell_{2}}^{2}\right) +\left(1-\eta\right)\left\|x\right\|_{\ell_{2}}^{2} \\
        & = \mathcal{K}\left(x\right)+\left(1-\eta\right)\left\|x\right\|_{\ell_{2}}^{2},
    \end{align*}
    where $\mathcal{K}\left(x\right):=\left\|x\right\|_{\ell_{1}}^{2}- \left\|x\right\|_{\ell_{2}}^{2}$. Thus, we observe that 
    \begin{align}\label{equ3.7}
        \mathcal{R}_{\eta}\left(x\right)-\mathcal{R}_{\eta}\left(x^{\dagger}\right)
        & = \mathcal{K}\left(x\right)-\mathcal{K}\left(x^{\dagger}\right) +\left(1-\eta\right)\left(\left\|x\right\|_{\ell_{2}}^{2}- \left\|x^{\dagger}\right\|_{\ell_{2}}^{2}\right) \nonumber\\
        & = \mathcal{K}\left(x\right)-\mathcal{K}\left(x^{\dagger}\right) +\left(1-\eta\right)\left\|x-x^{\dagger}\right\|_{\ell_{2}}^{2} +2\left(1-\eta\right)\left<x^{\dagger}, x-x^{\dagger}\right>.  
    \end{align}
    From the definition of $\mathcal{K}\left(x\right)$, we can express
    \begin{equation*}
        \mathcal{K}\left(x\right)-\mathcal{K}\left(x^{\dagger}\right) = \left\|x\right\|_{\ell_{1}}^{2}-\left\|x\right\|_{\ell_{2}}^{2} - \left\|x^{\dagger}\right\|_{\ell_{1}}^{2}+ \left\|x^{\dagger}\right\|_{\ell_{2}}^{2}.
    \end{equation*}
    With the definition of the index set $I\left(x^{\dagger}\right)$ in \eqref{equ2.2}, we can derive that
    \begin{align*}
      \left\|x\right\|_{\ell_{1}}^{2} & = \left(\sum_{i\in I\left(x^{\dagger}\right)} \left|x_{i}\right| + \sum_{i\notin I\left(x^{\dagger}\right)}\left|x_{i}\right| \right)^{2} \geq \left(\sum_{i\in I\left(x^{\dagger}\right)} \left|x_{i}\right| \right)^{2} + \sum_{i\notin I\left(x^{\dagger}\right)}\left|x_{i}\right|^{2}, \\
      -\left\|x\right\|_{\ell_{2}}^{2} & = -\sum_{i\in I\left(x^{\dagger}\right)} \left|x_{i}\right|^{2} - \sum_{i\notin I\left(x^{\dagger}\right)} \left|x_{i}\right|^{2}.
    \end{align*}
    Consequently, we obtain
    \begin{equation*}
        \mathcal{K}\left(x\right)-\mathcal{K}\left(x^{\dagger}\right) \geq \left(\sum_{i\in I\left(x^{\dagger}\right)}\left|x_{i}\right|\right)^{2} -\left(\sum_{i\in I\left(x^{\dagger}\right)} \left|x_{i}^{\dagger}\right|\right)^{2} - \sum_{i\in I\left(x^{\dagger}\right)} \left(\left|x_{i}\right|^{2} - \left|x_{i}^{\dagger}\right|^{2}\right),
    \end{equation*}
    which can be rewritten as
    \begin{align}\label{equ3.8}
        \mathcal{K}\left(x\right)-\mathcal{K}\left(x^{\dagger}\right) 
        & \geq \left(\sum_{i\in I\left(x^{\dagger}\right)}\left|x_{i}\right| + \sum_{i\in I\left(x^{\dagger}\right)} \left|x_{i}^{\dagger}\right|\right)\left(\sum_{i\in I\left(x^{\dagger}\right)}\left(\left|x_{i}\right|-\left|x_{i}^{\dagger}\right |\right)\right) \nonumber\\
        & \quad -\sum_{i\in I\left(x^{\dagger}\right)} \left(\left|x_{i}\right|+\left|x_{i}^{\dagger}\right|\right) \left(\left|x_{i}\right|-\left|x_{i}^{\dagger}\right|\right).
    \end{align}
    By the definition of $M_{1}$, we find
    \begin{equation*}
        \sum_{i\in I\left(x^{\dagger}\right)}\left|x_{i}\right| + \sum_{i\in I\left(x^{\dagger}\right)} \left|x_{i}^{\dagger}\right| \leq M_{1} + \left\|x^{\dagger}\right\|_{\ell_{1}}.
    \end{equation*}
    Thus, we conclude that
    \begin{equation}\label{equ3.9}
        \left|x_{i}\right| + \left|x_{i}^{\dagger}\right| \leq M_{1} + \left\|x^{\dagger}\right\|_{\ell_{1}}, \quad \forall i\in I\left(x^{\dagger}\right).
    \end{equation}
    Meanwhile, we observe that
    \begin{equation}\label{equ3.10}
        0<\left\|x^{\dagger}\right\|_{\ell_{1}}\leq \sum_{i\in I\left(x^{\dagger}\right)}\left|x_{i}\right| + \sum_{i\in I\left(x^{\dagger}\right)} \left|x_{i}^{\dagger}\right|.
    \end{equation}
    By combining inequalities \eqref{equ3.8}, \eqref{equ3.9} and \eqref{equ3.10}, we derive that
    \begin{equation*}
        \mathcal{K}\left(x\right)-\mathcal{K}\left(x^{\dagger}\right) \geq -\left\|x^{\dagger}\right\|_{\ell_{1}} \sum_{i\in I\left(x^{\dagger}\right)}\left|x_{i}-x_{i}^{\dagger}\right| - \left(M_{1}+\left\|x^{\dagger}\right\|_{\ell_{1}}\right) \sum_{i\in I\left(x^{\dagger}\right)}\left|x_{i}-x_{i}^{\dagger}\right|,
    \end{equation*}
    i.e.,
    \begin{equation}\label{equ3.11}
        \mathcal{K}\left(x\right)-\mathcal{K}\left(x^{\dagger}\right) \geq -\left(M_{1}+2\left\|x^{\dagger}\right\|_{\ell_{1}}\right) \sum_{i\in I\left(x^{\dagger}\right)}\left|x_{i}-x_{i}^{\dagger}\right|.
    \end{equation}
    A combination of \eqref{equ3.7} and \eqref{equ3.11} leads to the following inequality
    \begin{align}\label{equ3.12}
        \left(1-\eta\right)\left\|x-x^{\dagger}\right\|_{\ell_{2}}^{2}
        & \leq \mathcal{R}_{\eta}\left(x\right)-\mathcal{R}_{\eta}\left(x^{\dagger}\right) + \left(M_{1}+2\left\|x^{\dagger}\right\|_{\ell_{1}}\right) \sum_{i\in I\left(x^{\dagger}\right)}\left|x_{i}-x_{i}^{\dagger}\right| \nonumber\\
        & \quad -2\left(1-\eta\right)\left<x^{\dagger}, x-x^{\dagger}\right>.
    \end{align}
    By referring to Assumption \ref{ass3.1} (1), we find that
    \begin{equation*}
        \left|x_{i}-x^{\dagger}_{i}\right|=\left|\left<e_{i}, x-x^{\dagger}\right> \right| =\left|\left<\omega_{i}, F'\left(x^{\dagger}\right) \left(x-x^{\dagger}\right)\right>\right|\leq \max_{i\in I\left(x^{\dagger}\right)}\left\|\omega_{i}\right\|_{Y} \left\|F'\left(x^{\dagger}\right)\left(x-x^{\dagger}\right)\right\|_{Y}.
    \end{equation*}
    Thus, we can conclude that
    \begin{equation}\label{equ3.13}
        \sum_{i\in I\left(x^{\dagger}\right)}\left|x_{i}-x_{i}^{\dagger}\right| \leq \left|I\left(x^{\dagger}\right)\right|\max_{i\in I\left(x^{\dagger}\right)} \left\|\omega_{i}\right\|_{Y} \left\|F'\left(x^{\dagger}\right)\left(x-x^{\dagger}\right)\right\|_{Y},
    \end{equation} 
    where $\left|I\left(x^{\dagger}\right)\right|$ denotes the cardinality of the index set $I\left(x^{\dagger}\right)$. Additionally, according to Lemma \ref{lem3.3}, we can ascertain that
    \begin{equation}\label{equ3.14}
        \left|\left<x^{\dagger}, x-x^{\dagger}\right>\right| = \left|\left<\omega^{\dagger}, F'\left(x\right)\left(x-x^{\dagger}\right) \right>\right| \leq \left\|\omega^{\dagger}\right\|_{Y} \left\|F'\left(x^{\dagger}\right)\left(x-x^{\dagger}\right)\right\|_{Y}.
    \end{equation}
    A combination of \eqref{equ3.4}, \eqref{equ3.12}, \eqref{equ3.13} and \eqref{equ3.14} implies that
    \begin{align*}
        & \quad \left(1-\eta\right)\left\|x-x^{\dagger}\right\|_{\ell_{2}}^{2} \\
        & \leq \mathcal{R}_{\eta}\left(x\right)-\mathcal{R}_{\eta}\left(x^{\dagger}\right) + \left(\left(M_{1}+2\left\|x^{\dagger}\right\|_{\ell_{1}}\right) \left|I\left(x^{\dagger}\right)\right|\max_{i\in I\left(x^{\dagger}\right)}\left\|\omega_{i}\right\|_{Y} + 2\left(1-\eta\right)\left\|\omega^{\dagger}\right\|_{Y}\right) \\
        & \quad \cdot\left(\frac{\gamma}{2}\left\|x-x^{\dagger}\right\|_{\ell_{2}}^{2}+ \left\|F\left(x\right)-F\left(x^{\dagger}\right)\right\|_{Y}\right),
    \end{align*}
    i.e.,
    \begin{equation*}
        \left\|x-x^{\dagger}\right\|_{\ell_{2}}^{2} \leq \frac{1}{1-\frac{\gamma\left(c_{1}-c_{2}\eta\right)}{2\left(1-\eta\right)}} \left[\frac{1}{1-\eta}\left(\mathcal{R}_{\eta}\left(x\right)- \mathcal{R}_{\eta}\left(x^{\dagger}\right)\right)+ \frac{c_{1}-c_{2}\eta}{1-\eta} \left\|F\left(x\right)-F\left(x^{\dagger}\right)\right\|_{Y}\right],
    \end{equation*}
    where $c_{1}$ and $c_{2}$ are defined by \eqref{equ3.5} and \eqref{equ3.6}.
\end{proof}

\end{document}
